\documentclass[11pt,a4paper]{article}
\usepackage{patient_robot_instructions}

\title{%
  \vspace{-1.5cm}
  {\Large\bfseries Patient Instructions: Physical AI Oncology Trials}\\[0.3em]
  {\normalsize Instructional Sheets for 10 Robot Types — v1.9.1}
}
\author{%
  Kevin Kawchak\thanks{CEO, ChemicalQDevice. ORCID: \url{https://orcid.org/0009-0007-5457-8667}. Email: kevink@chemicalqdevice.com}\\
  \small DOI: \href{https://doi.org/10.5281/zenodo.18810541}{\texttt{10.5281/zenodo.18810541}}\\[0.2em]
  \small March 1, 2026
}
\date{}

\begin{document}

% ══════════════════════════════════════════════════════════════
% Each page is a self-contained instructional sheet for one robot type.
% Pages 5 and 6 feature pediatric patients.
% Robot types are ordered by frequency of use in physical AI oncology trials.
% Images are sourced from patients/images/ (numbered 1--10).
% ══════════════════════════════════════════════════════════════

% ──────────────────────────────────────────────────────────────
% PAGE 1: Surgical Robots
% ──────────────────────────────────────────────────────────────
\newpage
\thispagestyle{fancy}
\section*{Patient Instructions: Physical AI Oncology Trials - Surgical Robots}
\vspace{-0.3em}
\noindent\rule{\textwidth}{0.4pt}

\begin{center}
\vspace{0.3em}
\includegraphics[width=0.65\textwidth,keepaspectratio]{images/1.png}
\end{center}

\vspace{-0.3em}
\noindent\dashrule{\textwidth}
\begin{center}
{\large\bfseries Surgical Robots}
\end{center}
\noindent\rule{\textwidth}{0.4pt}

\vspace{0.3em}
This surgical robot assists your surgeon in removing a mediastinal (thymus) tumor through small chest incisions using precise multi-arm instruments.

\begin{enumerate}[leftmargin=1.5em,itemsep=2pt]
  \item Lie face-up on the operating table with your arms secured at your sides and remain still while anesthesia begins over about 2 minutes.
  \item You will be under general anesthesia for the full 90--180 minute procedure; the robot operates through 8--12\,mm chest ports while the surgeon controls it from a console.
  \item After the robot arms retract, you are moved to recovery and awaken in 15--30 minutes; keep your upper chest relaxed and follow the nurse's breathing checks.
\end{enumerate}

\vfill
\noindent\rule{\textwidth}{0.4pt}\\
{\scriptsize Source: Intuitive Surgical \hfill \textit{For Demonstration Purposes Only}}

% ──────────────────────────────────────────────────────────────
% PAGE 2: Cobots
% ──────────────────────────────────────────────────────────────
\newpage
\thispagestyle{fancy}
\section*{Patient Instructions: Physical AI Oncology Trials - Cobots}
\vspace{-0.3em}
\noindent\rule{\textwidth}{0.4pt}

\begin{center}
\vspace{0.3em}
\includegraphics[width=0.65\textwidth,keepaspectratio]{images/2.png}
\end{center}

\vspace{-0.3em}
\noindent\dashrule{\textwidth}
\begin{center}
{\large\bfseries Cobots (Collaborative Robots)}
\end{center}
\noindent\rule{\textwidth}{0.4pt}

\vspace{0.3em}
This collaborative robot arm stabilizes your forearm to assist a needle biopsy of a forearm soft-tissue sarcoma with steady positioning.

\begin{enumerate}[leftmargin=1.5em,itemsep=2pt]
  \item Sit in the exam chair, rest your forearm on the padded support palm-down, and keep your fingers relaxed for about 1 minute.
  \item The cobot moves at a maximum of 250\,mm/s with force-limited contact; hold still for 10--25 minutes as it repositions 2--3 times for sampling.
  \item The cobot retracts to its home position; remain seated for 2 minutes, then receive a bandage on your forearm before leaving.
\end{enumerate}

\vfill
\noindent\rule{\textwidth}{0.4pt}\\
{\scriptsize Source: Franka Robotics \hfill \textit{For Demonstration Purposes Only}}


% ──────────────────────────────────────────────────────────────
% PAGE 3: Radiotherapy Patient-Positioning Robots
% ──────────────────────────────────────────────────────────────
\newpage
\thispagestyle{fancy}
\section*{Patient Instructions: Physical AI Oncology Trials - RT Positioning Robots}
\vspace{-0.3em}
\noindent\rule{\textwidth}{0.4pt}

\begin{center}
\vspace{0.3em}
\includegraphics[width=0.65\textwidth,keepaspectratio]{images/3.png}
\end{center}

\vspace{-0.3em}
\noindent\dashrule{\textwidth}
\begin{center}
{\large\bfseries Radiotherapy Patient-Positioning Robots}
\end{center}
\noindent\rule{\textwidth}{0.4pt}

\vspace{0.3em}
This robotic treatment couch precisely positions your head with six degrees of freedom to align radiation for your brain tumor treatment.

\begin{enumerate}[leftmargin=1.5em,itemsep=2pt]
  \item Lie on the couch face-up with arms at your sides; allow the head mask to be fitted and secured over 2--3 minutes.
  \item The couch shifts in 1--3\,mm increments to align you; breathe normally and keep your head completely still during the 15--30 minute session.
  \item When the beam stops, the couch lowers to exit height; wait for the mask to be removed, then sit up slowly using both hands for support.
\end{enumerate}

\vfill
\noindent\rule{\textwidth}{0.4pt}\\
{\scriptsize Source: Accuray \hfill \textit{For Demonstration Purposes Only}}


% ──────────────────────────────────────────────────────────────
% PAGE 4: Robotic Needle-Placement Systems
% ──────────────────────────────────────────────────────────────
\newpage
\thispagestyle{fancy}
\section*{Patient Instructions: Physical AI Oncology Trials - Needle-Placement Robots}
\vspace{-0.3em}
\noindent\rule{\textwidth}{0.4pt}

\begin{center}
\vspace{0.3em}
\includegraphics[width=0.65\textwidth,keepaspectratio]{images/4.png}
\end{center}

\vspace{-0.3em}
\noindent\dashrule{\textwidth}
\begin{center}
{\large\bfseries Robotic Needle-Placement Systems}
\end{center}
\noindent\rule{\textwidth}{0.4pt}

\vspace{0.3em}
This robotic needle-guidance system uses CT imaging to place a biopsy needle with high accuracy into a cheek-area (parotid) tumor.

\begin{enumerate}[leftmargin=1.5em,itemsep=2pt]
  \item Lie on the imaging table face-up with your head on the pillow, keep your jaw relaxed, and hold still while the cheek is numbed for 1--2 minutes.
  \item The robot aligns and advances the needle over 1--3 minutes; do not talk or swallow when prompted during the 20--45 minute total procedure.
  \item After the needle is withdrawn, gentle pressure is applied for 3--5 minutes; remain lying down for 15--30 minutes while staff monitor swelling and pain.
\end{enumerate}

\vfill
\noindent\rule{\textwidth}{0.4pt}\\
{\scriptsize Source: ISO 15223-1 \hfill \textit{For Demonstration Purposes Only}}


% ──────────────────────────────────────────────────────────────
% PAGE 5: Social Companion Robots (PEDIATRIC)
% ──────────────────────────────────────────────────────────────
\newpage
\thispagestyle{fancy}
\section*{Patient Instructions: Physical AI Oncology Trials - Companion Robots}
\vspace{-0.3em}
\noindent\rule{\textwidth}{0.4pt}

\begin{center}
\vspace{0.3em}
\includegraphics[width=0.65\textwidth,keepaspectratio]{images/5.png}
\end{center}

\vspace{-0.3em}
\noindent\dashrule{\textwidth}
\begin{center}
{\large\bfseries Social Companion Robots}
\end{center}
\noindent\rule{\textwidth}{0.4pt}

\vspace{0.3em}
This friendly companion robot will talk and play interactive games with you to help you feel comfortable before your leukemia treatment today.

\begin{enumerate}[leftmargin=1.5em,itemsep=2pt]
  \item Walk into the room and wave or say ``Hello'' to the robot; it can offer a gentle high-five with its soft hand if you feel ready.
  \item Stand near the robot and play a short game for 10--20 minutes; follow its calm-breathing prompt while waiting for your visit.
  \item The robot says goodbye and shows a virtual sticker; give it a high-five or wave, then walk to the door where staff are waiting.
\end{enumerate}

\vfill
\noindent\rule{\textwidth}{0.4pt}\\
{\scriptsize Source: SoftBank Robotics \hfill \textit{For Demonstration Purposes Only}}


% ──────────────────────────────────────────────────────────────
% PAGE 6: Humanoids (PEDIATRIC)
% ──────────────────────────────────────────────────────────────
\newpage
\thispagestyle{fancy}
\section*{Patient Instructions: Physical AI Oncology Trials - Humanoids}
\vspace{-0.3em}
\noindent\rule{\textwidth}{0.4pt}

\begin{center}
\vspace{0.3em}
\includegraphics[width=0.65\textwidth,keepaspectratio]{images/6.png}
\end{center}

\vspace{-0.3em}
\noindent\dashrule{\textwidth}
\begin{center}
{\large\bfseries Humanoids}
\end{center}
\noindent\rule{\textwidth}{0.4pt}

\vspace{0.3em}
This humanoid robot practices gentle hand and balance exercises with you to help you prepare for osteosarcoma physical therapy sessions.

\begin{enumerate}[leftmargin=1.5em,itemsep=2pt]
  \item Approach the kneeling robot when it waves; hold its hands if comfortable and stand still together for about 1 minute.
  \item Follow its slow count-in/count-out breathing and posture cues for 15--25 minutes; keep your feet planted and shoulders relaxed.
  \item The robot releases your hands, waves, and stops; walk back to the waiting area where your parent or guardian is waiting.
\end{enumerate}

\vfill
\noindent\rule{\textwidth}{0.4pt}\\
{\scriptsize Source: Boston Dynamics \hfill \textit{For Demonstration Purposes Only}}


% ──────────────────────────────────────────────────────────────
% PAGE 7: Radiotherapy Motion-Management / Tracking Robots
% ──────────────────────────────────────────────────────────────
\newpage
\thispagestyle{fancy}
\section*{Patient Instructions: Physical AI Oncology Trials - RT Motion-Tracking Robots}
\vspace{-0.3em}
\noindent\rule{\textwidth}{0.4pt}

\begin{center}
\vspace{0.3em}
\includegraphics[width=0.65\textwidth,keepaspectratio]{images/7.png}
\end{center}

\vspace{-0.3em}
\noindent\dashrule{\textwidth}
\begin{center}
{\large\bfseries Radiotherapy Motion-Management / Tracking Robots}
\end{center}
\noindent\rule{\textwidth}{0.4pt}

\vspace{0.3em}
This motion-tracking system monitors your breathing so the radiation beam activates only when your lung tumor is aligned in the chest.

\begin{enumerate}[leftmargin=1.5em,itemsep=2pt]
  \item Lie on the treatment couch with arms at your sides; a reflective marker block is placed on your upper chest as the system calibrates in 1--2 minutes.
  \item Follow the breathing guide for 10--20 minutes; stay still between breaths so tracking stays within the 2--3\,mm tolerance.
  \item After the final beam, hold still while the marker block is removed from your chest, then sit up slowly using both hands on the couch edge.
\end{enumerate}

\vfill
\noindent\rule{\textwidth}{0.4pt}\\
{\scriptsize Source: Varian Medical \hfill \textit{For Demonstration Purposes Only}}

% ──────────────────────────────────────────────────────────────
% PAGE 8: Imaging Assistant Robots
% ──────────────────────────────────────────────────────────────
\newpage
\thispagestyle{fancy}
\section*{Patient Instructions: Physical AI Oncology Trials - Imaging Robots}
\vspace{-0.3em}
\noindent\rule{\textwidth}{0.4pt}

\begin{center}
\vspace{0.3em}
\includegraphics[width=0.65\textwidth,keepaspectratio]{images/8.png}
\end{center}

\vspace{-0.3em}
\noindent\dashrule{\textwidth}
\begin{center}
{\large\bfseries Imaging Assistant Robots}
\end{center}
\noindent\rule{\textwidth}{0.4pt}

\vspace{0.3em}
This robotic ultrasound system scans just under your ribs on the upper abdomen to measure a liver tumor and guide treatment planning.

\begin{enumerate}[leftmargin=1.5em,itemsep=2pt]
  \item Lie on the exam bed face-up with shoulders relaxed; keep both hands resting on your blanket while gel is applied over about 1 minute.
  \item The robotic probe applies 1--3\,N of steady pressure under the ribcage and scans for 10--20 minutes; breathe normally as it compensates for motion.
  \item The probe lifts away and gel is wiped with a warm cloth; you may sit up immediately---no recovery period is needed.
\end{enumerate}

\vfill
\noindent\rule{\textwidth}{0.4pt}\\
{\scriptsize Source: ISO 20417 \hfill \textit{For Demonstration Purposes Only}}

% ──────────────────────────────────────────────────────────────
% PAGE 9: Steerable Needle / Needle-Steering Robots
% ──────────────────────────────────────────────────────────────
\newpage
\thispagestyle{fancy}
\section*{Patient Instructions: Physical AI Oncology Trials - Steerable Needle Robots}
\vspace{-0.3em}
\noindent\rule{\textwidth}{0.4pt}

\begin{center}
\vspace{0.3em}
\includegraphics[width=0.65\textwidth,keepaspectratio]{images/9.png}
\end{center}

\vspace{-0.3em}
\noindent\dashrule{\textwidth}
\begin{center}
{\large\bfseries Steerable Needle / Needle-Steering Robots}
\end{center}
\noindent\rule{\textwidth}{0.4pt}

\vspace{0.3em}
This needle-steering robot guides a flexible needle through the upper abdomen to reach your liver tumor for targeted ablation treatment.

\begin{enumerate}[leftmargin=1.5em,itemsep=2pt]
  \item Lie face-up with your arms resting at your sides and hold still while local anesthetic is injected (brief 5-second sting).
  \item The needle advances while imaging confirms position every 5--10 seconds; remain motionless for the 30--60 minute procedure and breathe shallowly.
  \item The needle is withdrawn; apply gentle pressure to the gauze for 5--10 minutes, then rest on the table for 30--60 minutes of monitoring.
\end{enumerate}

\vfill
\noindent\rule{\textwidth}{0.4pt}\\
{\scriptsize Source: ISO 7010 \hfill \textit{For Demonstration Purposes Only}}

% ──────────────────────────────────────────────────────────────
% PAGE 10: Rehabilitation Exoskeletons / Robotic Gait Trainers
% ──────────────────────────────────────────────────────────────
\newpage
\thispagestyle{fancy}
\section*{Patient Instructions: Physical AI Oncology Trials - Rehab Exoskeletons}
\vspace{-0.3em}
\noindent\rule{\textwidth}{0.4pt}

\begin{center}
\vspace{0.3em}
\includegraphics[width=0.65\textwidth,keepaspectratio]{images/10.png}
\end{center}

\vspace{-0.3em}
\noindent\dashrule{\textwidth}
\begin{center}
{\large\bfseries Rehabilitation Exoskeletons / Robotic Gait Trainers}
\end{center}
\noindent\rule{\textwidth}{0.4pt}

\vspace{0.3em}
This lower-body exoskeleton supports your legs and guides walking practice to rebuild strength after femur osteosarcoma surgery.

\begin{enumerate}[leftmargin=1.5em,itemsep=2pt]
  \item Stand between the parallel bars, grip both rails, and hold steady while the waist, thigh, shank, and foot straps are secured over 2--3 minutes.
  \item Follow the step cues and walk at 0.2--0.5\,m/s for 15--30 minutes; say ``Pause'' to stop at any time while keeping both hands on the bars.
  \item The exoskeleton returns you to a standing rest position; keep gripping the rails as straps are removed (feet, shank, thigh, waist), then sit and rest for 5 minutes.
\end{enumerate}

\vfill
\noindent\rule{\textwidth}{0.4pt}\\
{\scriptsize Source: Ekso Bionics \hfill \textit{For Demonstration Purposes Only}}


\end{document}